L'\textbf{ingegneria del software} è una disciplina ingegneristica riguardante ogni aspetto relativo allo sviluppo del software professionale, dalla specifica del sistema fino alla manutenzione dopo che è entrato in uso. Aiuta a ottenere prodotti con la qualità richiesta entro i limiti di tempo e costo, accettando anche compromessi. Si tratta di un approccio sistematico e organizzato, orientato in base al caso di studio.\par 
Esistono infatti degli insiemi di attività, detti \textbf{processi di sviluppo}, atti alla creazione di un prodotto software. Si dividono in due macrocategorie: \textbf{plan-driven}, dove lo sviluppo è pianificato in modo dettagliato e le fasi sono seguite secondo un piano prestabilito, e \textbf{processi agili}, che garantiscono un notevole adattamento ai cambiamenti tramite un frequente confronto con il cliente. Nonostante la diversità, i processi di sviluppo condividono quattro attività fondamentali: \begin{itemize}
	\item \textbf{Specifiche del software}: Clienti e ingegneri definiscono funzionalità e vincoli del software.
	\item \textbf{Sviluppo del software}: Progettazione e implementazione del sistema.
	\item \textbf{Convalida del software}: Testing del software per garantire che rispetti le richieste del cliente.
	\item \textbf{Evoluzione del software}: Il prodotto viene modificato per soddisfare eventuali cambiamenti dei requisiti del cliente e del mercato.
\end{itemize}
\noindent Non esiste il processo perfetto e universalmente applicabile ad ogni caso di studio, ragion per cui nei prossimi capitoli studieremo diverse tecniche, le quali si potranno scegliere in base ai quattro aspetti più influenti per lo sviluppo: \begin{itemize}
	\item \textbf{Diversità}: Necessità che il software si adatti a diverse piattaforme o tecnologie.
	\item \textbf{Variabilità}: È verosimile che con il cambiare del mercato, i committenti modifichino a loro volta i requisiti. La sfida della variabilità consiste nell'essere agili abbastanza da sviluppare e modificare il software per adattarlo a queste casistiche.
	\item \textbf{Fiducia e protezione}: Essendo il software radicato in ogni aspetto della vita è necessario che funzioni correttamente e non sia vulnerabile ad attacchi di malintenzionati.
	\item \textbf{Scala}: Il software deve essere sviluppato per sistemi di dimensioni molto diverse, dai piccoli dispositivi fino a grandi sistemi distribuiti e servizi cloud.
\end{itemize}
\noindent Come ci si può aspettare, lo sviluppo di software professionale di grandi dimensioni è infattibile da un singolo individuo: per questa ragione, in condizioni ideali, le aziende creano una task-force per la gestione del lavoro, composta dalle seguenti figure: \begin{itemize}
	\item \textbf{Analista}: Definisce i requisiti decisi insieme al committente.
	\item \textbf{Designer}: Si occupa della strutturazione di design ad alto e basso livello.
	\item \textbf{Programmatore}: Scrive il codice.
	\item \textbf{Tester}: Verifica la correttezza del codice.
	\item \textbf{Mantenitore}: Si occupa della manutenzione del software dopo il rilascio, correggendo errori, migliorando il sistema e adattandolo a nuovi requisiti o tecnologie.
	\item \textbf{Project Manager}: Responsabile della pianificazione e controllo di un progetto. Gestisce inoltre il personale, stima i costi e cerca di far rispettare budget e tempi. Idealmente saprebbe giostrarsi bene in ognuno dei campi precedenti.
\end{itemize}
\noindent Non sempre lo sviluppo del software va come pianificato; possono presentarsi problemi di vario tipo, come per esempio sforare il budget iniziale o i tempi di consegna. Ciò può accadere per una moltitudine di motivi, ma da vari report emerge che le cause principali derivano da una comunicazione non esaustiva tra \textbf{stakeholder} e team operativo, dalla difficoltà nel comprendere il dominio di lavoro, oppure da problemi organizzativi nel team.\par
Come vedremo nei seguenti capitoli, sarà di fondamentale importanza definire appropriatamente i requisiti e useremo vari sottoprocessi per "estrarre" quante più richieste possibile dai committenti. In sistemi normali, la loro validazione avviene tra analista e stakeholder, ma esistono software il cui funzionamento è fondamentale che sia garantito. Parliamo in questo caso di \textbf{sistemi critici}, come macchinari utilizzati in ambito medico, per i quali il software deve rispettare standard rigorosi ed essere certificato da enti di validazione indipendenti.